\documentclass[12pt]{article}
\usepackage[utf8]{inputenc}
\usepackage{amsmath}
\usepackage{amssymb}
\usepackage{geometry}
\usepackage{xcolor}
\usepackage{tikz}
\usetikzlibrary{shapes,arrows}

\geometry{a4paper, margin=1in}

\title{\textbf{\Lsquare 勾股定理的几何证明}}
\author{\textit{DeepSeek-R1 + Manim}}
\date{\today}

\begin{document}

\maketitle

\section{引言}

勾股定理是平面几何中最基本且最重要的定理之一。它描述了直角三角形三条边之间的关系：

\begin{equation}
\label{eq:pythagorean}
a^2 + b^2 = c^2
\end{equation}

其中：
\begin{itemize}
    \item $a$ 和 $b$ 是直角三角形的两条直角边
    \item $c$ 是斜边（直角所对的边）
\end{itemize}

\section{几何证明：矩形对角线}

考虑一个宽度为 $a$、高度为 $b$ 的矩形。我们可以通过矩形的对角线来验证勾股定理。

\subsection{矩形构造}

设矩形 $ABCD$ 满足：
\begin{itemize}
    \item $AB = a = 3$ （宽度）
    \item $AD = b = 4$ （高度）
    \item $AC = c$ （对角线长度，待求）
\end{itemize}

\subsection{对角线长度计算}

根据勾股定理，对角线长度为：

\begin{align}
c &= \sqrt{a^2 + b^2} \\
  &= \sqrt{3^2 + 4^2} \\
  &= \sqrt{9 + 16} \\
  &= \sqrt{25} \\
  &= 5
\end{align}

\subsection{数值验证}

\begin{equation}
3^2 + 4^2 = 9 + 16 = 25 = 5^2
\end{equation}

\section{可视化}

使用 Manim 动画引擎，我们可以直观地展示这个几何关系：

\begin{itemize}
    \item \textcolor{blue}{蓝色矩形}：表示 $3 \times 4$ 的矩形
    \item \textcolor{red}{红色对角线}：连接矩形的对角顶点，长度为 5
    \item \textcolor{yellow}{黄色边框}：突出显示勾股定理公式
\end{itemize}

\section{结论}

通过矩形的对角线，我们验证了勾股定理的正确性。对于 $3 \times 4$ 的矩形：

\begin{center}
\Large \textcolor{green}{$\boxed{3^2 + 4^2 = 5^2}$}
\end{center}

\end{document}
